\chapter{Ανάλυση και Σχεδίαση του Συστήματος}
\label{chap:design}

\section{Σχεδιαστικός στόχος}
\label{sec:design-goal}
Το αντικείμενο της σχεδίασης είναι ένα σύστημα επίλυσης του κύβου 3×3×3. Πρόκειται για ένα σύνολο επιλυτών (\eng{solvers}) με διαφορετικές εγγυήσεις, που καλύπτουν το φάσμα από αποδεδειγμένα βέλτιστες έως γρήγορες πρακτικές λύσεις. Μαζί τους περιλαμβάνεται η υποδομή αναζήτησης και πινάκων που τους τροφοδοτεί. Ο πρώτος σχεδιαστικός στόχος είναι επομένως λειτουργικός: ένας αξιόπιστος πυρήνας μοντέλου κύβου και αναζήτησης, πάνω στον οποίο στηρίζονται οι αλγόριθμοι του Κεφαλαίου~\ref{chap:algorithms}.

Ο δεύτερος στόχος αφορά την ακεραιότητα των ισχυρισμών. Στη βέλτιστη επίλυση είναι εύκολο να ξεφύγουν αβάσιμοι ισχυρισμοί: είναι εύκολο να γραφτεί ότι ένας επιλυτής είναι βέλτιστος επειδή βρήκε σύντομη λύση. Είναι πιο δύσκολο, αλλά ακαδημαϊκά απαραίτητο, να δείξει κανείς ποια διαδικασία απέδειξε το μήκος της λύσης, ποιο αρχείο αποθηκεύει το αποτέλεσμα και ποιος έλεγχος το επαλήθευσε. Για τον λόγο αυτό η σχεδίαση εξασφαλίζει ότι κάθε αποτέλεσμα είναι ιχνηλάσιμο, από τις εντολές παραγωγής έως τους πίνακες του Κεφαλαίου~\ref{chap:experiments}. Αν ένα αποτέλεσμα δεν μπορεί να αναπαραχθεί, τότε ακόμη και ένας σωστός αριθμός δεν στηρίζει αξιόπιστα τον ισχυρισμό.

\section{Επίπεδα αρχιτεκτονικής}
\label{sec:design-architecture}
Η αρχιτεκτονική οργανώνεται σε επίπεδα, όπως συνοψίζει το Σχήμα~\ref{fig:architecture}. Κάθε επίπεδο προσφέρει στο επόμενο μια περιορισμένη διεπαφή, ώστε οι ευθύνες να παραμένουν διακριτές και κάθε αλλαγή να έχει προβλέψιμο εύρος.

\begin{figure}[htbp]
\centering
% Αρχιτεκτονική του συστήματος σε επίπεδα (περιεχόμενο σχήματος· το περιβάλλον
% figure και η λεζάντα ορίζονται στο Κεφάλαιο 5). Τα ονόματα αρθρωμάτων
% αντιστοιχούν στο πραγματικό δέντρο πηγαίου κώδικα: src/rubik_optimal/
% (cli.py, oracle.py, verify.py, solvers/, tables/), native/ και results/.
\begin{tikzpicture}[
  font=\footnotesize,
  band/.style={draw=black!60, rounded corners=3pt, fill=black!4},
  bandtitle/.style={anchor=north west, font=\scriptsize\bfseries, text=black!70, inner sep=2pt},
  comp/.style={draw=black!75, rounded corners=1.5pt, fill=white, align=center,
               font=\scriptsize, inner sep=3pt, minimum height=1.05cm},
  flow/.style={-{Stealth[length=2.4mm]}, semithick},
  resflow/.style={-{Stealth[length=2.4mm]}, semithick, black!55},
  lbl/.style={font=\tiny, align=center}
]

% Επίπεδο 1: Διεπαφή γραμμής εντολών
\draw[band] (0,14.3) rectangle (12.4,15.9);
\node[bandtitle] at (0.12,15.82) {Διεπαφή γραμμής εντολών (CLI)};
\node[comp, minimum width=10.8cm] at (6.2,14.85)
  {{\ttfamily rubik-optimal} ({\ttfamily src/rubik\_optimal/cli.py}): οκτώ υποεντολές\\
   {\tiny\ttfamily solve · distance · oracle · benchmark · verify · scramble · facelets · tables}};

% Επίπεδο 2: Επιλυτές Python
\draw[band] (0,11.3) rectangle (12.4,13.5);
\node[bandtitle] at (0.12,13.42) {Επιλυτές Python ({\ttfamily rubik\_optimal.solvers})};
\node[comp, minimum width=2.85cm] at (1.65,12.2)
  {Korf / IDA*\\ {\tiny\ttfamily korf.py}};
\node[comp, minimum width=2.85cm] at (4.68,12.2)
  {Kociemba\\ {\tiny\ttfamily kociemba.py}\\ {\tiny\ttfamily kociemba\_optimal.py}\\
   {\tiny προσαρμογέας {\ttfamily kociemba} (PyPI)}};
\node[comp, minimum width=2.85cm] at (7.71,12.2)
  {Thistlethwaite\\ {\tiny\ttfamily thistlethwaite.py}};
\node[comp, minimum width=2.85cm] at (10.74,12.2)
  {Κύβος 2×2×2 (\eng{Pocket})\\ {\tiny\ttfamily pocket\_cube.py · pocket/}};

% Στρώμα μαντείου (γέφυρα επιπέδων 2-3)
\node[comp, fill=black!8, minimum width=7.4cm] at (8.4,10.05)
  {Στρώμα μαντείου: {\ttfamily oracle.py}: {\ttfamily FastOptimalOracle}\\
   {\tiny παραμένουσα διεργασία (\eng{resident}) και λειτουργία δέσμης (\eng{batch})}};

% Επίπεδο 3: Εγγενείς μηχανές C/C++
\draw[band] (1.6,6.85) rectangle (12.4,8.85);
\node[bandtitle] at (1.72,8.77) {Εγγενείς μηχανές C/C++ ({\ttfamily native/})};
\node[comp, minimum width=2.5cm] at (3.01,7.6)
  {Εγγενής IDA*\\ {\tiny\ttfamily optimal\_solver}};
\node[comp, minimum width=2.5cm] at (5.67,7.6)
  {Παραγωγή PDB\\ {\tiny\ttfamily corner\_pdb}\\ {\tiny\ttfamily edge\_pdb}};
\node[comp, minimum width=2.5cm] at (8.33,7.6)
  {Απόδειξη φάσης 2\\ {\tiny\ttfamily kociemba\_phase2\_probe}};
\node[comp, minimum width=2.5cm] at (10.99,7.6)
  {Μηχανή H48\\ {\tiny\ttfamily h48\_backend}\\ {\tiny ενσωματωμένο {\ttfamily nissy-core}}};

% Επίπεδο 4: Πίνακες
\draw[band] (0,4.05) rectangle (12.4,6.05);
\node[bandtitle] at (0.12,5.97) {Πίνακες ({\ttfamily data/generated/})};
\node[comp, minimum width=3.7cm] at (2.17,4.8)
  {PDB γωνιών + 8 PDB έξι ακμών\\
   {\tiny\ttfamily *\_corner\_state\_pdb.bin}\\ {\tiny\ttfamily *\_edge\_subset\_*\_pdb.bin}};
\node[comp, minimum width=3.7cm] at (6.19,4.8)
  {Πίνακες αποστάσεων\\ συμπλόκων\\
   {\tiny Thistlethwaite · φάση 1 Kociemba}};
\node[comp, minimum width=3.7cm] at (10.21,4.8)
  {Πίνακας αποκοπής H48\\ {\tiny\ttfamily h48h7.bin}};

% Επίπεδο 5: Επαλήθευση και αποτελέσματα
\draw[band] (0,1.05) rectangle (12.4,3.25);
\node[bandtitle] at (0.12,3.17) {Επαλήθευση και αποτελέσματα};
\node[comp, minimum width=3.7cm] at (2.17,2.0)
  {Ελεγκτής λύσεων\\ {\tiny\ttfamily verify.py}\\ {\tiny\ttfamily scripts/verify\_results.py}};
\node[comp, minimum width=3.7cm] at (6.19,2.0)
  {Αρχεία αποτελεσμάτων\\ {\tiny\ttfamily results/raw $\to$ results/processed}};
\node[comp, minimum width=3.7cm] at (10.21,2.0)
  {Εξωτερικός διασταυρωτικός\\ έλεγχος: Nissy 2.0.8\\ {\tiny\ttfamily solvers/nissy\_external.py}};

% Ροές κλήσεων (προς τα κάτω)
\draw[flow] (6.2,14.3) -- (6.2,13.5)
  node[lbl, right=2pt, midway] {κλήσεις επιλυτών και μαντείου};
\draw[flow] (8.4,11.3) -- (8.4,10.5)
  node[lbl, right=2pt, midway] {ερωτήματα ακριβούς απόστασης /\\ βέλτιστης λύσης};
\draw[flow] (8.4,9.6) -- (8.4,8.85)
  node[lbl, right=2pt, midway] {παραμένουσα εγγενής διεργασία};
\draw[flow] (3.0,11.3) -- (3.0,8.85)
  node[lbl, left=2pt, midway, align=right]
  {περιτυλίγματα\\ {\ttfamily optimal\_native.py}\\ {\ttfamily h48\_native.py}};
\draw[flow] (0.7,11.3) -- (0.7,6.05);
\node[lbl, rotate=90] at (0.42,8.68) {απευθείας φόρτωση πινάκων};
\draw[flow] (6.19,6.85) -- (6.19,6.05)
  node[lbl, right=2pt, midway] {παραγωγή και φόρτωση πινάκων};

% Ροή αποτελεσμάτων (δεξιός διάδρομος)
\draw[resflow] (12.4,12.2) -- (13.25,12.2) -- (13.25,2.0) -- (12.4,2.0);
\draw[semithick, black!55] (12.4,7.6) -- (13.25,7.6);
\node[lbl, rotate=90, text=black!60] at (13.6,7.1)
  {ροή εγγραφών αποτελεσμάτων (JSONL)};

\end{tikzpicture}

\caption{Η αρχιτεκτονική του συστήματος σε επίπεδα: διεπαφή γραμμής εντολών, επιλυτές Python, εγγενείς (\eng{native}) μηχανές C/C++, στρώμα πινάκων αποκοπής και πινάκων προτύπων, στρώμα μαντείου (\eng{oracle}), και αλυσίδα επαλήθευσης και αποτελεσμάτων με εξωτερικό διασταυρωτικό έλεγχο.}
\label{fig:architecture}
\end{figure}

Το πρώτο επίπεδο είναι το μοντέλο του κύβου μέσα στο πακέτο Python \code{rubik_optimal}. Περιλαμβάνει την αναπαράσταση κατάστασης σε επίπεδο κυβιδίων (\eng{cubies}), τις κινήσεις, τον έλεγχο φυσικής εγκυρότητας και τις μετατροπές από και προς τη συμβολοσειρά ψηφίδων (\eng{facelet string}), στα αρχεία \code{cube.py}, \code{moves.py} και \code{validity.py}. Το στρώμα συμμετριών βρίσκεται στο \code{symmetry.py}. Οι συμβάσεις του μοντέλου ορίζονται στο Κεφάλαιο~\ref{chap:model}.

Το δεύτερο επίπεδο είναι η υποδομή αναζήτησης και ευρετικών στο υποπακέτο \code{search}: πλήρης BFS μικρού βάθους, παραμετρική IDA* και κοινές συναρτήσεις ευρετικής. Σε αυτό προστίθεται το στρώμα συντεταγμένων και πινάκων στα υποπακέτα \code{coordinates} και \code{tables}. Εκεί παράγονται και διαβάζονται οι πίνακες κινήσεων, οι πίνακες αποκοπής (\eng{pruning tables}) και οι πίνακες προτύπων (\eng{pattern databases}, PDB). Ο συνδυασμός IDA* με παραδεκτούς πίνακες προτύπων ακολουθεί τη σχεδίαση του Korf \cite{korf1997pattern}.

Το τρίτο επίπεδο είναι οι επιλυτές Python στο υποπακέτο \code{solvers}. Σε αυτό ανήκουν η υλοποίηση Korf/IDA* (\code{korf.py}), η υλοποίηση του αλγορίθμου δύο φάσεων κατά Kociemba (\code{kociemba.py}) \cite{kociembaTwoPhase}, η αποδεικτικά βέλτιστη παραλλαγή της ίδιας σχεδίασης (\code{kociemba_optimal.py}), η τετραφασική υλοποίηση Thistlethwaite (\code{thistlethwaite.py}) και ο επιλυτής του κύβου 2×2×2 (\eng{Pocket Cube}), που στηρίζεται στο υποπακέτο \code{pocket}. Στο ίδιο επίπεδο ανήκει ο προσαρμογέας (\eng{adapter}) προς το εξωτερικό πακέτο \code{kociemba} του PyPI \cite{muodovKociembaPkg}, ο οποίος χρησιμεύει ως εξωτερική βάση σύγκρισης για πρακτικές λύσεις.

Το τέταρτο επίπεδο είναι οι μηχανές C/C++, οι οποίες παράγουν τα κύρια ακριβή αποτελέσματα της εργασίας. Στον κατάλογο \code{native/} βρίσκονται πέντε συστατικά: ο παραγωγός του γωνιακού πίνακα προτύπων (\code{native/corner_pdb}), ο παραγωγός των πινάκων προτύπων έξι ακμών (\code{native/edge_pdb}), ο επιλυτής IDA* πλήρους κύβου σε C (\code{native/optimal_solver}), η μηχανή απόδειξης βελτιστότητας πάνω στη σχεδίαση δύο φάσεων (\code{native/kociemba_phase2_probe}) και η ενσωματωμένη (\eng{vendored}) μηχανή H48 του \code{nissy-core} (\code{native/h48_backend}) \cite{trontoNissyCoreH48}. Οι μηχανές αυτές οδηγούνται από την Python μέσω των περιβλημάτων \code{solvers/optimal_native.py} και \code{solvers/h48_native.py}.

Το πέμπτο επίπεδο είναι το στρώμα μαντείου. Το αρχείο \code{oracle.py} παρέχει τη διεπαφή \code{FastOptimalOracle}, η οποία απαντά σε ερωτήματα ακριβούς απόστασης και βέλτιστης λύσης για αυθαίρετες έγκυρες καταστάσεις 3×3×3. Χρησιμοποιεί από προεπιλογή τον πίνακα \code{h48h7}, διατηρεί μια παραμένουσα διεργασία C για τα επαναλαμβανόμενα ερωτήματα και κρύβει από τον κώδικα που το καλεί τις λεπτομέρειες των μηχανών C/C++.

Το έκτο επίπεδο είναι η αλυσίδα πειραμάτων, επαλήθευσης και αποτελεσμάτων. Το \code{benchmark.py} παράγει τις μετρήσεις, το \code{verify.py} επαναλαμβάνει κάθε λύση στο μοντέλο κυβιδίων, και τα αποτελέσματα αποθηκεύονται σε αρχεία με σταθερό σχήμα (Παράρτημα~\ref{app:schema}). Στο ίδιο επίπεδο ανήκει ο εξωτερικός διασταυρωτικός έλεγχος: το \code{solvers/nissy_external.py} συνδέει το σύστημα με το εξωτερικό εκτελέσιμο \code{nissy}, ώστε τα ακριβή αποτελέσματα μεγάλου βάθους να επιβεβαιώνονται και από ανεξάρτητη υλοποίηση (Κεφάλαιο~\ref{chap:validation}). Πάνω από όλα τα επίπεδα βρίσκεται η διεπαφή γραμμής εντολών (\code{cli.py}), η οποία δίνει ενιαία πρόσβαση στους επιλυτές, στην αναγνώριση απόστασης (\eng{distance recognition}) και στο μαντείο (Παράρτημα~\ref{app:cli}).

Ο διαχωρισμός σε επίπεδα κρατά τις ευθύνες χωριστές. Αν αλλάξει η αναπαράσταση κατάστασης, επηρεάζονται οι κινήσεις και οι κωδικοποιητές συντεταγμένων, αλλά όχι η μορφή των πινάκων αποτελεσμάτων. Αν αλλάξει το σώμα μετρήσεων (\eng{benchmark corpus}), αλλάζουν τα ακατέργαστα αποτελέσματα και οι παραγόμενοι πίνακες, αλλά όχι το μοντέλο κυβιδίων. Οι επιλυτές δεν γράφουν απευθείας στο κείμενο της εργασίας και τα σενάρια παραγωγής δεν παρακάμπτουν τον έλεγχο λύσεων. Το σύστημα δεν κυνηγά πάση θυσία τις επιδόσεις λογισμικού επιπέδου παραγωγής (\eng{production})· προτεραιότητα έχει η σαφήνεια του συμβολαίου κάθε επιπέδου.

\section{Ροή δεδομένων}
\label{sec:design-dataflow}
Η βασική ροή δεδομένων αρχίζει από ντετερμινιστικές περιπτώσεις ελέγχου. Το \code{benchmark.py} δημιουργεί τις περιπτώσεις με καθορισμένο σπόρο (\eng{seed}) και προφίλ εκτέλεσης (\eng{profile}). Για κάθε περίπτωση, το \code{benchmark.py} παράγει την κατάσταση κυβιδίων από το ανακάτεμα (\eng{scramble}), τη μετατρέπει σε συμβολοσειρά ψηφίδων και την αποθηκεύει στην αντίστοιχη εγγραφή αποτελέσματος. Στη συνέχεια εκτελούνται οι επιλυτές και η αναγνώριση απόστασης. Κάθε εκτέλεση συμπληρώνει την εγγραφή με τη λύση, ετικέτα κατάστασης, χρόνο εκτέλεσης, πλήθος κόμβων και ένδειξη επαλήθευσης.

Οι ακατέργαστες εγγραφές γράφονται σε μορφή JSONL, με μία γραμμή ανά πειραματικό αποτέλεσμα. Η μορφή αυτή επιτρέπει εύκολη επιθεώρηση και απλή επεξεργασία. Από τις ακατέργαστες εγγραφές παράγονται συγκεντρωτικές συνόψεις και αρχεία CSV, και από αυτές οι πίνακες LaTeX και τα δεδομένα των σχημάτων της εργασίας. Η ροή είναι μονόδρομη: οι πίνακες LaTeX δεν αποτελούν ποτέ πηγή νέων αποτελεσμάτων — είναι το τελικό παραγόμενο αρχείο, όχι δεδομένο εισόδου.

Η ροή των πινάκων συντεταγμένων είναι ανάλογη, αλλά ακολουθεί δικό της δρόμο. Οι προδιαγραφές συντεταγμένων ορίζουν πεδίο τιμών, συναρτήσεις κωδικοποίησης και αποκωδικοποίησης, λυμένη τιμή και περιγραφή. Η κοινή γεννήτρια παράγει πίνακα κινήσεων και πίνακα αποκοπής για κάθε προδιαγραφή, αποθηκεύει τα αρχεία των πινάκων και γράφει μανιφέστο (\eng{manifest}). Το μανιφέστο τροφοδοτεί τον αντίστοιχο παραγόμενο πίνακα της εργασίας, ώστε το κείμενο να παρουσιάζει μεγέθη και πλήθη εγγραφών χωρίς πρόσβαση στους ίδιους τους μεγάλους πίνακες. Οι επιλυτές φορτώνουν ή κατασκευάζουν τους πίνακες αποκοπής μόνο μέσω κοινών συναρτήσεων ευρετικής.

Οι πίνακες των μηχανών C/C++ ακολουθούν το ίδιο πρότυπο σε μεγαλύτερη κλίμακα. Τα σενάρια παραγωγής μεταγλωττίζουν τις μηχανές C/C++, παράγουν τους δυαδικούς πίνακες προτύπων και τους πίνακες H48, υπολογίζουν αθροίσματα ελέγχου SHA-256 και γράφουν μανιφέστο στα επεξεργασμένα αποτελέσματα. Οι επιλυτές και το μαντείο καταναλώνουν τους πίνακες αυτούς μόνο μέσω του μανιφέστου και των περιβλημάτων τους. Έτσι κάθε χρήση πίνακα μπορεί να ανιχνευθεί πίσω στο συγκεκριμένο αρχείο και στο άθροισμα ελέγχου του.

\section{Προφίλ εκτέλεσης και καταστάσεις αναπαραγωγιμότητας}
\label{sec:design-profiles}
Το σύστημα διαχωρίζει τρία προφίλ εκτέλεσης: \code{quick}, \code{thesis} και \code{stress}. Το προφίλ \code{quick} προορίζεται για ανάπτυξη και γρήγορους ελέγχους, με μικρότερες περιπτώσεις και μικρότερα χρονικά όρια. Το προφίλ \code{thesis} είναι η κύρια πηγή των αριθμών του κειμένου. Το προφίλ \code{stress} προορίζεται για βαθύτερη διερεύνηση πέρα από το κύριο σώμα μετρήσεων.

Ο διαχωρισμός αυτός αποτρέπει δύο σφάλματα. Το πρώτο θα ήταν να χρησιμοποιηθούν πρόχειρα αποτελέσματα ανάπτυξης ως τεκμήρια της εργασίας. Το δεύτερο θα ήταν να απαιτείται εκτενής εκτέλεση μετρήσεων για κάθε μικρή αλλαγή του κειμένου. Με τα προφίλ, η καθημερινή ανάπτυξη παραμένει γρήγορη, ενώ το τελικό κείμενο βασίζεται στο προφίλ \code{thesis}. Τα ονόματα των αρχείων αποτελεσμάτων περιέχουν το προφίλ και τον σπόρο, ώστε να φαίνεται από πού προέρχεται κάθε αρχείο αποτελεσμάτων.

\section{Στρατηγική ορθότητας}
\label{sec:design-correctness}
Η ορθότητα εξασφαλίζεται σε πολλά επίπεδα. Στο χαμηλότερο επίπεδο υπάρχουν δοκιμές μονάδας για τις κινήσεις, τον αναλυτή ακολουθιών κινήσεων και τη φυσική εγκυρότητα. Στο επόμενο επίπεδο υπάρχουν δοκιμές για τις αμφίδρομες μετατροπές συντεταγμένων και τις μεταβάσεις των πινάκων, και σε ακόμη υψηλότερο επίπεδο δοκιμές επιλυτών για περιπτώσεις μικρού βάθους. Κανένα επίπεδο δεν υποκαθιστά τα άλλα: μια δοκιμή μονάδας που περνά δεν αποδεικνύει ότι οι παραγόμενοι πίνακες είναι ενημερωμένοι, και μια μεταγλώττιση του κειμένου που πετυχαίνει δεν αποδεικνύει ότι οι ισχυρισμοί είναι σωστοί. Το σύστημα χρειάζεται πολλαπλούς ελέγχους επειδή κάθε έλεγχος καλύπτει διαφορετικό είδος αστοχίας.

Στο πειραματικό επίπεδο, στην αρχιτεκτονική ο ελεγκτής αποτελεσμάτων (\eng{verifier}) δεν παίρνει τίποτα για δεδομένο από κανέναν επιλυτή: κάθε καταγεγραμμένη λύση επαναλαμβάνεται στο μοντέλο κυβιδίων πριν θεωρηθεί έγκυρη. Το σενάριο \code{scripts/verify_results.py} εφαρμόζει τον έλεγχο αυτόν στις ακατέργαστες εγγραφές. Επιπλέον επικυρώνει: τους πίνακες προτύπων και τους πίνακες αποκοπής· τα μεταδεδομένα των πινάκων H48· την απόδειξη βελτιστότητας των λύσεων 3×3×3 από τις μηχανές C/C++· τη σύνοψη του κύβου 2×2×2· και την παρουσία των παραγόμενων αρχείων πινάκων της εργασίας. Αν ο έλεγχος περάσει, η λύση είναι έγκυρη — δεν σημαίνει όμως ότι είναι και βέλτιστη. Ο μηχανισμός της επαλήθευσης περιγράφεται στο Κεφάλαιο~\ref{chap:implementation}, ενώ οι εγγυήσεις και τα όρια της σημασιολογίας των ακριβών αποτελεσμάτων αναλύονται στο Κεφάλαιο~\ref{chap:validation}.

Η σχεδίαση των ετικετών κατάστασης είναι μέρος της ίδιας στρατηγικής. Το σχήμα αποτελεσμάτων προβλέπει, μεταξύ άλλων, τις εξής ετικέτες: \code{exact} (αποδεδειγμένα βέλτιστη λύση)· \code{non_exact} (επαληθευμένη αλλά όχι αποδεδειγμένα βέλτιστη λύση)· \code{lower_bound} (αποδεδειγμένο κάτω φράγμα χωρίς λύση)· και \code{timeout} (εξάντληση χρονικού ή υπολογιστικού ορίου). Ο πλήρης κατάλογος δίνεται στο Παράρτημα~\ref{app:schema}. Αν ένας επιλυτής επιστρέψει λύση χωρίς απόδειξη βελτιστότητας, το σύστημα έχει τρόπο να το δηλώσει. Αν μια κατάσταση είναι φυσικώς αδύνατη (μη έγκυρη), η αναγνώριση απόστασης δεν την αντιμετωπίζει σαν να ήταν απλώς μια δύσκολη περίπτωση. Η ίδια αυστηρότητα που είναι ενσωματωμένη στον μηχανισμό αποτυπώνεται και στο κείμενο της εργασίας.

\section{Στρατηγική παραγωγής πινάκων}
\label{sec:design-tables}
Οι πίνακες συντεταγμένων έχουν κεντρικό ρόλο στην εργασία, επειδή συνδέουν τη θεωρία των ευρετικών με την πρακτική αναζήτηση. Η παραγωγή τους γίνεται με κοινή γεννήτρια και όχι με χειροκίνητη αποθήκευση. Το σύστημα γράφει τόσο το πλήρες περιεχόμενο κάθε πίνακα όσο και τα μεταδεδομένα του, ώστε το κείμενο να τεκμηριώνει μεγέθη και πλήθη εγγραφών χωρίς πρόσβαση στα δυαδικά αρχεία.

Η σχεδίαση των προδιαγραφών επιτρέπει επέκταση. Μια νέα συντεταγμένη χρειάζεται κωδικοποίηση, αποκωδικοποίηση, μέγεθος πεδίου τιμών, λυμένη τιμή και περιγραφή. Αν προστεθεί στο σύνολο των παραγόμενων προδιαγραφών, παράγεται με τον ίδιο μηχανισμό. Αυτό δεν σημαίνει ότι κάθε νέα συντεταγμένη είναι αλγοριθμικά χρήσιμη ή αποδοτική· σημαίνει όμως ότι η υποδομή δεν περιορίζεται μόνο στις τρεις αρχικές προβολές.

Το ότι παρήχθησαν προβολές και για τη φάση 2 δείχνει την αξία αυτής της επέκτασης. Οι αρχικές προβολές CO/EO/UD-slice επαρκούν για τη φάση 1 και για γενικό κάτω φράγμα. Η δεύτερη φάση του αλγορίθμου δύο φάσεων χρειάζεται πληροφορία για τις μεταθέσεις μέσα στο περιορισμένο σύνολο κινήσεων \cite{kociembaTwoPhase}. Προστέθηκαν επομένως προβολές μετάθεσης γωνιών, μετάθεσης ακμών U/D και μετάθεσης ακμών της μεσαίας φέτας για τη φάση 2. Η ίδια λογική επέκτασης ακολουθείται και στους πίνακες προτύπων C/C++, κατά την προσέγγιση του Korf \cite{korf1997pattern}: ο γωνιακός πίνακας καλύπτει πλήρως τη γωνιακή προβολή και οι οκτώ πίνακες έξι ακμών προσθέτουν ουσιαστική πληροφορία για τις ακμές. Το ίδιο σχήμα θα μπορούσε να δεχθεί ισχυρότερες συνδυασμένες συντεταγμένες ή συντεταγμένες με κατάτμηση κόστους (\eng{cost partitioning}) σε μελλοντική εργασία.

\section{Σχεδίαση του σώματος μετρήσεων}
\label{sec:design-corpus}
Το σώμα μετρήσεων του προφίλ \code{thesis} επιλέχθηκε ώστε να είναι μικρό αλλά να επιτρέπει σαφή ερμηνεία των αποτελεσμάτων. Περιλαμβάνει τη λυμένη κατάσταση, ντετερμινιστικές περιπτώσεις μικρού βάθους και βαθύτερα τυχαία ανακατέματα. Κάθε κατηγορία απαντά σε διαφορετικό ερώτημα: η λυμένη κατάσταση ελέγχει την τετριμμένη συμπεριφορά, οι περιπτώσεις μικρού βάθους τη συνέπεια των ακριβών αποτελεσμάτων, και οι περιπτώσεις μεγαλύτερου βάθους αναδεικνύουν τις πρακτικές διαφορές των επιλυτών και τα όρια των αναζητήσεων.

Το σώμα αυτό δεν είναι αντιπροσωπευτικό δείγμα ολόκληρου του χώρου καταστάσεων του 3×3×3, και αυτό δηλώνεται ρητά. Η αξία του είναι η αναπαραγωγιμότητα και η ερμηνευσιμότητα. Μια μεγαλύτερη εμπειρική μελέτη θα απαιτούσε δείγμα πολλών εκατοντάδων ή χιλιάδων ανακατεμάτων, στατιστική ανάλυση των χρόνων εκτέλεσης και ενδεχομένως διαφορετικά χρονικά όρια. Η επέκταση αυτή συζητείται στο Κεφάλαιο~\ref{chap:conclusions}. Το τρέχον σώμα είναι κατάλληλο για να δείξει ότι οι υλοποιήσεις λειτουργούν και πώς διαφέρουν οι εγγυήσεις τους.

\section{Αναπαραγώγιμη παραγωγή πινάκων αποτελεσμάτων}
\label{sec:design-artifacts}
Στο κείμενο της εργασίας κανένας αριθμός δεν αντιγράφεται με το χέρι. Οι πίνακες αποτελεσμάτων και τα δεδομένα των σχημάτων παράγονται από τα επεξεργασμένα αρχεία και εισάγονται στο LaTeX ως παραγόμενα αρχεία. Έτσι η ανανέωση μιας μέτρησης ενημερώνει το κείμενο χωρίς χειροκίνητη επεξεργασία πινάκων, αποκλείοντας τα λάθη αντιγραφής. Η προσέγγιση απαιτεί πειθαρχία: μετά από κάθε αλλαγή αποτελεσμάτων πρέπει να αναπαράγονται οι πίνακες και να επαναλαμβάνονται οι έλεγχοι. Τον ρόλο αυτόν υποστηρίζει το σενάριο \code{scripts/thesis_audit.py}. Το σενάριο αυτό ελέγχει την έκταση του κειμένου, την παρουσία και τη χρονική συνέπεια των παραγόμενων αρχείων και την προέλευσή τους. Εντοπίζει επίσης φράσεις-ισχυρισμούς μέσα στο κείμενο που χρειάζονται έλεγχο από άνθρωπο. Οι δείκτες αυτοί δεν είναι αναγκαστικά λάθη· είναι σημεία όπου ο αυτόματος έλεγχος δεν μπορεί να αντικαταστήσει την κρίση του συγγραφέα. Η πλήρης διαδικασία αναπαραγωγής των αποτελεσμάτων περιγράφεται στο Παράρτημα~\ref{app:repro}.

\section{Διαχείριση περιορισμών}
\label{sec:design-limitations}
Οι περιορισμοί δεν μπαίνουν απλώς σε ένα κεφάλαιο στο τέλος — είναι ενσωματωμένοι στη σχεδίαση. Το σχήμα αποτελεσμάτων διαθέτει ετικέτες κατάστασης, ώστε οι αναζητήσεις που δεν κατέληξαν σε λύση να μην αποκρύπτονται. Παράλληλα, το κείμενο επαναλαμβάνει τα όρια εκεί όπου υπάρχει κίνδυνος παρερμηνείας.

Η επανάληψη αυτή είναι σκόπιμη. Σε μια εργασία για βέλτιστη επίλυση, ο αναγνώστης πρέπει να βλέπει σε κάθε σημείο αν ένα αποτέλεσμα είναι ακριβές ή πρακτικό. Αν η διάκριση εμφανιζόταν μόνο σε ένα κεφάλαιο περιορισμών, οι πίνακες και τα συμπεράσματα θα μπορούσαν να διαβαστούν λανθασμένα. Για τον λόγο αυτό η ορολογία των ετικετών \code{exact}, \code{non_exact}, \code{lower_bound} και \code{timeout} συνοδεύει τα αποτελέσματα από τη σχεδίαση έως τα πειράματα και τα συμπεράσματα.

\section{Επεκτασιμότητα}
\label{sec:design-extensibility}
Η αρχιτεκτονική αφήνει περιθώριο για πληρέστερη υλοποίηση χωρίς να ανατραπεί το υπάρχον σύστημα. Μπορούν να προστεθούν νέες συντεταγμένες, νέοι πίνακες αποκοπής, νέοι επιλυτές ή νέα προφίλ εκτέλεσης. Είναι κρίσιμο να ακολουθεί κάθε επέκταση το ίδιο συμβόλαιο: δοκιμές, παραγόμενα τεκμήρια, επαλήθευση λύσεων, παραγωγή πινάκων και ειλικρινή δήλωση κατάστασης. Αν προστεθεί ισχυρότερος πίνακας για τη φάση 2, πρέπει να αποτυπωθεί στο μανιφέστο. Αν μειωθούν οι μη ολοκληρωμένες αναζητήσεις, πρέπει να φανεί σε νέα σύνοψη μετρήσεων. Αν αλλάξει κάποιος ισχυρισμός, πρέπει να αλλάξει και το κείμενο.

Αντίστοιχα, η μετάβαση σε μεγαλύτερους πίνακες προτύπων θα απαιτούσε αλλαγή στον τρόπο αποθήκευσης, αλλά όχι στη βασική φιλοσοφία. Το μανιφέστο μπορεί να κρατά αθροίσματα ελέγχου και την κατάσταση των πηγών, οι επιλυτές να εξακολουθούν να χρησιμοποιούν κοινές συναρτήσεις ευρετικής, και η επαλήθευση να ελέγχει τις επιστρεφόμενες λύσεις. Το σημαντικό είναι να μην καταστεί αδιαφανές το ισχυρότερο παραγόμενο αρχείο: ακόμη και αν ένας πίνακας αποθηκεύεται σε δυαδική μορφή, πρέπει να συνοδεύεται από μανιφέστο και σενάριο παραγωγής.

\section{Σύνοψη της σχεδίασης}
\label{sec:design-summary}
Η σχεδίαση του συστήματος υπηρετεί την ακαδημαϊκή θέση της εργασίας: μια υλοποίηση είναι χρήσιμη μόνο όταν οι εγγυήσεις της είναι σαφείς. Στόχος δεν είναι να εντυπωσιάσει το σύστημα με αδιαφανείς αριθμούς. Στόχος είναι κάθε αριθμός να μπορεί να ανιχνευθεί και κάθε αποτέλεσμα επιλυτή να κατατάσσεται ανάλογα με την εγγύηση που προσφέρει. Η προσέγγιση αυτή είναι αυστηρότερη από μια απλή επίδειξη επίλυσης, αλλά ταιριάζει καλύτερα σε μια εργασία όπου το ζητούμενο είναι τεκμηριωμένη γνώση και όχι απλώς ένα σύστημα που δουλεύει.
